\section{Arduino}
\subsection{Introductions to Arduino and microcontrollers}
\subsubsection*{Microcontrollers}
Microcontrollers are complete computer systems with microprocessor (CPU), clock oscillator, memory and I/O peripherals integrated
on a single chip. Microcontrollers are used to get data from the environment through sensors, process it and control actuators
accordingly. Due to memory and processing power limitations, code for microcontrollers must be written in a very efficient way.

\subsubsection*{Arduino Uno}
Arduino Uno is a microcontroller board based on the ATmega328P microcontroller. It has 14 digital input/output pins and 6 analog
input pins and a 16 MHz quartz crystal. It has an AREF pin for analog reference voltage for ADC and a reset button. It has a USB
connection for programming and a power jack for external power supply. It has a 5V pin (400 mA) and a 3.3V pin (50 mA) connected
to the onboard voltage regulator and a Vin pin connected directly to the input power supply source.

The Atmega328P microcontroller has 32 KB of ISP flash memory for storing code, 2 KB of SRAM for storing variables and 1 KB of
EEPROM for storing non volatile data. It has 23 GPIO (general purpose input/output), 32 general purpose working registers,
6 channel 10-bit ADC, USART, SPI and I2C communication interfaces. It has three timers, internal and external interrupts and
a watchdog timer. It operates at a voltage between 1.8 and 5.5 V.

\subsubsection*{Arduino Mega}
Arduino Mega is a microcontroller board based on the ATmega2560 microcontroller. It has 54 digital input/output pins and 16
analog input pins and a 16 MHz quartz crystal. Due to the higher number of pins, it is used for more complex projects that
require more I/O peripherals.

\subsubsection*{Arduino Due}
Arduino Due is a microcontroller board based on the Atmel SAM3X8E ARM Cortex-M3 CPU. It has 54 digital input/output pins and 12
analog input pins and a 84 MHz quartz crystal. Due to the higher processing power, it is used for more complex projects that
require more processing power and more I/O peripherals. \textbf{It operates at a voltage of 3.3 V and it is not 5V tolerant.}

\subsubsection*{Arduino Nano 33 BLE Sense Rev2}
Arduino Nano 33 BLE Sense Rev2 is a microcontroller board based on the Nordic Semiconductor nRF52840 ARM Cortex-M4 CPU in the
same form factor as the Arduino Nano. It has 14 digital input/output pins and 8 analog input pins and a 64 MHz quartz crystal.
It has a Bluetooth Low Energy (BLE) module and a wide range of sensors, including an accelerometer, a gyroscope, a magnetometer,
a temperature sensor, a humidity sensor, a pressure sensor and a microphone and has an RGB LED. \textbf{It operates at a voltage
of 3.3 V and it is not 5V tolerant.}

\subsubsection*{Other boards}
\begin{itemize}[topsep=-1pt, itemsep=-1pt]
	\item Arduino Uno Wifi Rev2: same form factor as the Uno but with an ATmega4809 microcontroller and a WiFi module with
	integrated TCP/IP protocol stack.
	\item Arduino STAR - Otto: equiped with a powerful ARM processor to control power demanding peripherals like LCD, Audio, Camera, etc.
	\item Arduino Alicepad: board for wearables and clothing
	\item Arduino Nano: same microcontroller as the Uno but in a smaller form factor
	\item Arduino Giga R1 Wifi: board with an ARM Cortex-M7 microcontroller and a WiFi module, designed for high performance
	applications like robotics, AI and machine learning
	\item Arduino Portenta H7: board with an ARM Cortex-M7 microcontroller and a WiFi module, designed for high performance
	applications like robotics, AI and machine learning
	\item Shields: passive boards that can be plugged on top of an Arduino board to extend its functionality
\end{itemize}

\subsection{Programming on Arduino - pin modes, interrupts and timers}
\subsubsection*{Bootloader}
The bootloader is a small program stored in a private section of the flash memory that allows to program the microcontroller
through a simple serial communication. On the board, there is also an USB-to-serial converter that allows to connect the
a computer to the microcontroller via USB and use it as a serial communication interface for programming and debugging.

Usually the bootloader is pre-flashed on the microcontroller by the manufacturer. If the bootloader is not present or needs
to be updated, it must be flashed on the microcontroller via the SPI communication pins (MOSI, MISO, SCK) connected to an
external programmer.

\subsubsection*{Program structure}
An arduino program is called a sketch and it is written in C/C++. It has two main functions:
\begin{itemize}
	\item \texttt{setup()}: executed only once, when the microcontroller is powered on or reset; it is used to initialize
	variables, object instances of peripherals, pin modes and other settings
	\item \texttt{loop()}: executed repeatedly after the \texttt{setup()} function; it is used to implement the main logic
	of the program
\end{itemize}

\noindent
Considerations about the program structure:
\begin{itemize}
	\item global variables are preferrred to local variables, because they are stored in the flash memory and do not consume
	SRAM, which is a very limited resource on microcontrollers
\end{itemize}

\subsubsection*{Pin modes and pin functions}
Arduino pins can be configured in different modes and can perform different functions.
\begin{itemize}
	\item digital output: pin configured as a digital output, used to write digital signals (LOW or HIGH) through
	\texttt{digitalWrite(PIN, [HIGH,LOW])} function
	\item digital input: pin configured as a digital input, used to read digital signals (HIGH or LOW) through
	\texttt{digitalRead(PIN)} function
	\item analog output: pin configured as a digital PWM output, used to write PWM signals with custom duty cycle in
	range 0-255 through \texttt{analogWrite(PIN, value)} function, PWM frequency can vary depending on board and pin used
	\item analog input: pin configured as an analog input, used to read analog signals through the ADC and get
	a digital value in range 0-1023 (for 10-bit ADC) through \texttt{analogRead(PIN)} function
\end{itemize}

\noindent
To configure how the pin behaves, the \texttt{pinMode(PIN, MODE)} function is used in the \texttt{setup()} function, where
\texttt{MODE} can be one of the following:
\begin{itemize}
	\item \texttt{OUTPUT}: configures the pin as a digital output
	\item \texttt{INPUT}: configures the pin as an analog/digital input
	\item \texttt{INPUT\_PULLUP}: configures the pin as a digital input with an internal pull-up resistor, in this mode the pin
	is connected to \(V_{CC}\) through a \(\approx 22 \; k\Omega\) pull-up resistor, so the pin reads HIGH when it is not connected to
	ground and LOW when it is connected to ground.
\end{itemize}

\subsubsection*{Analog readings and ADC configuration}
The resolution of the ADC can be changed by calling the \texttt{analogReadResolution(bits)} function in the \texttt{setup()}
function. By default, the resolution is 10 bits (0-1023), but it can be increased to 12 bits (0-4095) on some boards like the
Arduino Due and the Arduino Nano 33 BLE Sense Rev2.

It is possible to set an external high-voltage reference for the ADC by connecting that reference voltage to the AREF pin and
calling the \texttt{analogReference(EXTERNAL)} function in the \texttt{setup()} function. By default, the reference voltage is
the supply voltage (5V or 3.3V depending on the board).

The resolution step of the ADC can be calculated as follows:
\[\text{ADC step} = \frac{V_{REF}}{2^{\text{bits}}} \qquad \text{by default} \begin{array}{l l}
	\text{Arduino Uno: } & \displaystyle \text{ADC step} = \frac{5 \; V}{2^{10 \; \text{bits}}} = 4.88 \; mV \\[7pt]
	\text{Arduino Nano 33 BLE: } & \displaystyle \text{ADC step} = \frac{3.3 \; V}{2^{12 \; \text{bits}}} = 0.806 \; mV
\end{array}\]

\subsubsection*{The \texttt{millis}, \texttt{micros} and \texttt{delay} functions}
These three functions are used to manage time in Arduino sketches and use internal timers of the microcontroller.
\begin{itemize}
	\item \texttt{millis()}: returns the number of milliseconds since the microcontroller was powered on or reset;
	it returns an unsigned 32-bit integer, so it overflows after approximately 50 days
	\item \texttt{micros()}: returns the number of microseconds since the microcontroller was powered on or reset;
	it returns an unsigned 32-bit integer, so it overflows after approximately 70 minutes; due to limited clock speed
	on some boards, the resolution of this function can be 4 microseconds (16 MHz clock)
	\item \texttt{delay(ms)}: blocks the execution of the program for a specified number of milliseconds
\end{itemize}

\noindent
Use of the \texttt{millis()} and \texttt{micros()} functions is preferred to the use of the \texttt{delay()} function,
because it does not block the execution of the program and allows to implement non-blocking delays and timeouts, which
are essential for implementing real-time applications.

\subsubsection*{Interrupts}
An interrupt is a signal that can be generated by an external event (for example a button press) or by an internal event
(for example a timer overflow) that interrupts the normal execution of the program and executes a special function called
interrupt service routine (ISR). Interrupts are used to handle events that require immediate attention and cannot be handled
in the loop function of the program, for example user interactions.

\noindent
Interrupt service routines has the following characteristics:
\begin{itemize}
	\item they must be as short and fast as possible
	\item only one interrupt can be executed at a time so if multiple interrupts are generated at the same time, they are
	executed in order of priority
	\item only \texttt{delayMicroseconds()} function works inside an ISR, while \texttt{delay()} and \texttt{millis()}
	functions do not work because they rely on interrupts to update the internal timers
	\item they do not accept parameters and do not return values, so they must use \texttt{volatile} global variables to
	communicate with the main program. The keyword \texttt{volatile} is used to prevent the compiler from optimizing operations
	on those variables and causing race conditions.
\end{itemize}

\noindent
To attach an interrupt to a pin, the \texttt{attachInterrupt(digitalPinToInterrupt(PIN), ISR, MODE)} function is used in the
\texttt{setup()} function where:
\begin{itemize}
	\item \texttt{PIN} is the pin number to which the interrupt is attached (only supported pins can be used)
	\item \texttt{ISR} is the name of the interrupt service routine function
	\item \texttt{MODE} is the condition that triggers the interrupt and can be one of the following:
	\begin{itemize}[topsep=-1pt]
		\item \texttt{RISING}: triggers the interrupt when the pin goes from LOW to HIGH
		\item \texttt{FALLING}: triggers the interrupt when the pin goes from HIGH to LOW
		\item \texttt{CHANGE}: triggers the interrupt when the pin changes state (from LOW to HIGH or vice versa)
		\item \texttt{LOW}: triggers the interrupt when the pin is LOW (only for some boards)
		\item \texttt{HIGH}: triggers the interrupt when the pin is HIGH (only for some boards)
	\end{itemize}
\end{itemize}

\subsubsection*{Switch-case, \texttt{break} \texttt{continue} statements}
Switch-case statements are used to execute different blocks of code based on the value of a variable. They are more efficient
than if-else statements when there are many cases to handle, because they use a jump table instead of evaluating multiple
conditions.

The \texttt{break} and \texttt{continue} statements are used to control the flow of loops and switch-case statements. The
\texttt{break} statement is used to exit a loop or a switch-case statement, while the \texttt{continue} statement is used
to skip the current iteration of a loop and continue with the next iteration.

\newpage

\subsection{Simple circuits with Arduino}
\subsubsection*{Drive a LED directly connected to an Arduino pin}
To correctly bias a LED, a current limiting resistor must be used in series with the LED. The value of the resistor can be
calculated using Ohm's law as follows.

\begin{center}
	\begin{minipage}{0.5\textwidth}
		\centering \begin{circuitikz}[american, scale=0.8, transform shape]
			\circuitikzset{diodes/scale=0.8}
			\draw (0,0) node[left] {\(\begin{array}{c} V_{CC} \\ \text{\footnotesize Arduino pin} \end{array}\)} to[short, o-] ++(0.5,0) to[R, l=$R$] ++(2,0) to[leD, l=$LED$] ++(2, 0) node[ground] {};
		\end{circuitikz}
	\end{minipage}
	\begin{minipage}{0.45\textwidth}
		\[R = \frac{V_R}{I_R} = \frac{V_{CC} - V_{LED}}{I_{LED}}\]
	\end{minipage}
\end{center}

\subsubsection*{Drive high current into a LED with a BJT}
Arduino pins can source or sink a limited amount of current (4mA for Nano 33 BLE Sense Rev2). If a higher current is needed,
a BJT can be used as a switch to drive the LED. The base of the BJT is connected to the Arduino pin through a current limiting
resistor, the collector is connected to the LED with a current limiting resistor in series and the emitter is connected to ground.
The transistor is used as a switch in saturation regime, so the voltage drop across the collector-emitter junction is negligible
and the current through the LED is determined by the current limiting resistor and the supply voltage.

\begin{center}
	\begin{minipage}{0.5\textwidth}
		\centering \begin{circuitikz}[american, scale=0.8, transform shape]
			\circuitikzset{diodes/scale=0.8}
			\circuitikzset{transistors/scale=1.4}
			\draw (0,0) node[left] {\(V_{CC}\)} to[short, o-] ++(0.5,0) to[R, l=$R_{LED}$] ++(2,0) to[leD, l=$LED$] ++(2, 0) node[npn, anchor=collector, rotate=90] (npn) {};
			\draw (npn.emitter) -- ++(0.5, 0) node[ground] {};
			\draw (npn.base) to[R, l=$R_B$] ++(0,-1.5) to[short, -o] ++(0, -0.5) node[left] {\(\begin{array}{c} V_{CC} \\ \text{\footnotesize Arduino pin} \end{array}\)};

			\node [xshift=0.2cm, yshift=0.1cm] at (npn.B) (vplus) {$+$};
			\node [xshift=-0.2cm, yshift=-0.2cm] at (npn.E) (vminus) {$-$};
			\draw[-{Latex}] (vminus) to[out=-90, in=10] node[pos=0.3, fill=white, inner sep=1pt] {$V_{BE}$} (vplus);
		\end{circuitikz}
	\end{minipage}
	\begin{minipage}{0.45\textwidth}
		\[R_{LED} = \frac{V_{R_{LED}}}{I_{R_{LED}}} = \frac{V_{CC} - V_{LED}}{I_{LED}}\]
		\vspace{5pt}
		\[R_B = \frac{V_{R_B}}{I_{R_B}} = \frac{V_{CC} - V_{BE}}{I_B}\]
	\end{minipage}
\end{center}

\subsubsection*{Dimming an LED / RGB-LED with \texttt{analogWrite()} and PWM signal}
By using the \texttt{analogWrite(PIN, value)} function, it is possible to generate a PWM signal on a digital pin (with PWM support)
and control the brightness of a LED by changing the duty cycle of the PWM signal. Value can be in range 0-255, where 0 means 0\% duty
cycle (LED off) and 255 means 100\% duty cycle (LED fully on).

On the Arduino Nano 33 BLE Sense Rev2, there is a built-in RGB LED controlled by three virtual pins (\texttt{LED\_R},
\texttt{LED\_G}, \texttt{LED\_B}). By using the \texttt{analogWrite()} function on these pins, it is possible to control
the brightness of each color and create different colors by mixing the three primary colors.

\subsubsection*{Read a button state with \texttt{digitalRead()}, pull-up resistors and debouncing}
A button is a simple switch that can be used to connect a pin to ground when pressed. The main problem is that when the button
is not pressed, the pin is left floating and can read random values. To solve this problem, a pull-up resistor should be used
to connect the pin to \(V_{CC}\). In this way, when the button is not pressed, the pin reads HIGH and when the button is pressed,
the pin reads LOW.

\begin{center}
	\begin{circuitikz}[american, scale=0.8, transform shape]
		\draw (0,0) node[left] {\(V_{CC}\)} to[short, o-] ++(0.5,0) to[R, l=$R_{PU}$] ++(2,0) to[push button, l=\texttt{button}] ++(2, 0) node[ground] {};
		\draw (2.5, 0) to[short, -o] ++(0, 1) node[above] {\texttt{digital pin}};
	\end{circuitikz}
\end{center}

To choose the value of the pull-up resistor, it is necessary to consider the current that flows through the resistor when the
button is pressed (when the button is not pressed, the current is negligible because the pin is in high impedance state). The
suggested current is in the order of 0.1 - 0.2 mA, so the value of the pull-up resistor can be calculated as follows.
\[I_{R_{PU}} = \frac{V_{CC}}{R_{PU}} \qquad \text{common values are}: \begin{array}{l}
	20-22 \; k\Omega \text{ for } 5 \; V \text{ boards} \\[7pt]
	10-12 \; k\Omega \text{ for } 3.3 \; V \text{ boards}
\end{array}\]

On some boards like the Arduino Uno, there are internal pull-up resistors of 20 k\(\Omega\) that can be enabled by configuring
the pin mode as \texttt{INPUT\_PULLUP}. In this way the button can be connected directly between the pin and ground, without
the need of an external pull-up resistor.

Usually, when a button is pressed, the transition from HIGH to LOW is not clean and immediate, but there are some oscillations
called bouncing due to mechanical switch contacts oscillations. To solve this problem, a simple software debouncing technique
can be used, which consists in waiting for a short time (10 and 50 ms) between two consecutive readings of the button state.

Alternatively, a pull-down resistor can be used between the pin and ground and the button will be connected between the pin
and \(V_{CC}\). The button behaviour becomes active-HIGH, meaning the state will be HIGH when the button is pressed and LOW
when the button is not pressed. 

\subsubsection*{Read a potentiometer value with \texttt{analogRead()}}
It is possible to read the voltage value of a potentiometer wiper terminal by connecting it to an analog input pin. The voltage
at the wiper terminal will be in range \([0, V_{CC}]\) and can be read by an analog pin using the \texttt{analogRead()} function.
The value returned is in range \([0, \;1023]\) for a 10-bit ADC or in range \([0, \; 4095]\) for a 12-bit ADC and is linearly
proportional to the wiper voltage.

In some applications, it is necessary to scale the value read from the potentiometer to a different range (for example \([0, \; 255]\)
to control the brightness of an LED). To do this, it is possible to use the \texttt{map(x, in\_min, in\_max, out\_min, out\_max)}
function to map the value \(x\) from the input range \([\texttt{in\_min}, \; \texttt{in\_max}]\) to the output range
\([\texttt{out\_min}, \; \texttt{out\_max}]\) as follows:
\[\texttt{map(x, in\_min, in\_max, out\_min, out\_max)} \; = \; \frac{(\texttt{x} - \texttt{in\_min}) \times (\texttt{out\_max} - \texttt{out\_min})}{\texttt{in\_max} - \texttt{in\_min}} + \texttt{out\_min}\]

\subsubsection*{Read analog sensors with \texttt{analogRead()}}
Sometimes the output voltage of an analog sensor has a smaller range \([0, V_{sensor,max} < V_{CC}]\). To correctly read the
sensor value, it is necessary to scale the sensor output voltage to the full range of the ADC. This can be done in different ways:
\begin{itemize}
	\item using the \texttt{map()} function to map the value read from the sensor, but does not increase the resolution of the ADC
	because is done after the ADC conversion
	\item using an external op-amp to amplify the sensor output voltage to the full range of the ADC; attention must be paid to
	not exceed the maximum input voltage of the ADC
	\item using the AREF pin to set the reference voltage of the ADC to the maximum sensor output voltage, so the ADC will
	automatically scale the sensor output voltage to the full range of the ADC (this one is the preferred solution)
\end{itemize}

\subsubsection*{Read temperature with analog TMP36 temperature sensor}
TMP36 is an analog temperature sensor that outputs a voltage linearly proportional to the temperature. The output voltage has
a 0.5 V offset (0.5 V correspond to 0\(^\circ C\)) and a slope of 10 mV/\(^\circ C\).
To read the temperature with an Arduino, the output voltage of the TMP36 sensor can be read with an analog input pin and
converted into temperature using the following mapping formula considering (155 correspond to 0.5 V and 330 is the temperature
range for 3.3 V power supply).
\[T = \frac{V_{out} - 0.5 \; V}{10 \; mV/^\circ C} = (V_{out} - 0.5 \; V) \times 100 \; ^\circ C/V \quad \rightarrow \quad \texttt{temperature = (value-155) * 330 / 1023;}\]

\subsubsection*{Read humidity with analog HIH5030 humidity sensor}
HIH5030 is an analog humidity sensor that outputs a voltage linearly proportional to the relative humidity. The output voltage
can be read with an analog input pin and converted into relative humidity using one of the following mapping formulas (from the datasheet).
\[\begin{array}{c}
	RH = (V_{out} / V_{CC} - 0.1515) / 0.00636 \\[5pt]
	V_{out} = \text{measured\_value} \cdot V_{CC} / 1023
\end{array} \quad \rightarrow \quad \texttt{humidity = (value / 1023 - 0.1515) / 0.00636;}\]

\newpage

\subsection{Asynchronous Serial Protocol}
\subsubsection*{Introduction to asynchronous serial protocols}
Asyncronous Serial Protocol is a family of communication protocols that allow to exchange data between the microcontroller
and other devices using a serial communication interface. It has the following characteristics:
\begin{itemize}
	\item asynchronous: the data is transmitted without a clock signal, so the sender and the receiver must agree before the
	communication on the baud rate (bits per second) and the format of the data (start and stop bits, parity bit, etc.)
	\item serial: the data is transmitted serially, one bit at a time, not in parallel
\end{itemize}

\subsubsection*{UART - Universal Asynchronous Receiver-Transmitter protocol}
The UART protocol is the asynchronous serial protocol used for the communication between an Arduino board and other device.
It is a full-duplex protocol that uses two lines for communication: TX (transmit) and RX (receive). The TX line of one device is connected to the RX line
of the other device and vice versa. In idle state, lines are HIGH. The communication has the following format:
\begin{itemize}
	\item start bit: 1 bit LOW, indicates the start of a data frame
	\item data bits: bits transmitted in little endian order (least significant bit first)
	\item stop bit: 1 bit HIGH, indicate the end of a data frame
\end{itemize}

\noindent

\subsubsection*{Serial communication with Arduino}
In Arduino boards, the UART protocol is implemented on the digital pins 0 (RX) and 1 (TX) for communication with other devices
and on the USB port for communication with the computer via the USB-to-serial converter. The Arduino IDE provides a built-in
Serial Monitor or a Serial plotter to visualize the data sent from the Arduino to the computer through the USB port and to send
data from the computer to the Arduino.

In Arduino sketches, the UART protocol is accessed through the \texttt{Serial} class, which provides the following functions:
\begin{itemize}
	\item \texttt{Serial.begin(baud\_rate)}: initializes the serial communication with the specified baud rate, usually used
	in the \texttt{setup()} function, common baud rates are 300, 600, 1200, 2400, 4800, 9600, 14400, 19200, 28800, 38400, 57600,
	115200 bps 
	\item \texttt{Serial.print(data)}: sends data to the serial port as human-readable ASCII text
	\item \texttt{Serial.println(data)}: same as \texttt{Serial.print()}, plus append a newline character at the end
	\item \texttt{Serial.available()}: returns if there is data available to read from the serial port
	\item \texttt{Serial.read()}: reads a byte of data from the serial port; returns -1 if no data is available
\end{itemize}

\newpage

\subsection{Serial Peripheral Interface - SPI}
\subsubsection*{Introduction to SPI protocol}
Serial Peripheral Interface (SPI) is a synchronous serial communication protocol that uses a master-slave architecture to exchange
data between a single controller device (master) and one or more controlled peripheral devices (slaves).

\subsubsection*{SPI connection lines}
The SPI protocol uses at least four lines for communication:
\begin{itemize}
	\item SCLK (Serial Clock): generated by the master to synchronize the communication
	\item CIPO (Controller In Peripheral Out): transmit data from the peripheral to the controller (ex MISO)
	\item COPI (Controller Out Peripheral In): transmit data from the controller to the peripheral (ex MOSI)
	\item CS (Chip Select): line used to select a specific peripheral device (active LOW)
\end{itemize}
SCLK, CIPO and COPI lines are shared between all devices on the SPI bus, while each peripheral device has its own CS line
connected to the master.

\subsubsection*{SPI communication process}
The SPI communication woks as follows:
\begin{itemize}
	\item the master selects the slave device by pulling its CS line LOW
	\item the master generates the clock signal on the SCLK line
	\item on each clock edge, the master and the slave exchange one bit of data on the CIPO and COPI lines
	\item after the communication is finished, the master stops the clock signal and pulls the CS line HIGH
\end{itemize}

\subsection{Libraries in Arduino}
\subsubsection*{Introduction to Arduino/C/C++ libraries}
A library is a collection of functions and declarations that provide extra functionality which can be used in multiple
arduino sketches without the need to write the code from scratch. Libraries are organized in:
\begin{itemize}
	\item header file (.h): contains the declarations of the functions and variables provided by the library
	\item source file (.cpp): contains the implementation of the functions declared in the header file
\end{itemize}

\subsubsection*{Classes and objects in Arduino libraries}
Usually libraries contain classes that are user-defined data structures that encapsulate data and functions and are used to
manage peripherals and sensors or software tools in a more organized way. Classes are defined in the header file and implemented
in the source file. Classes are instantiated in the arduino sketch to create objects of that class, which can be used to call
the functions provided by the library and manage the peripheral or sensor associated with that class.
Classes have the following components:
\begin{itemize}
	\item class constructor: a special function that is called when an object of the class is instantiated, it initializes
	the object and its variables
	\item public functions: functions that can be called from outside the class, they are used to provide the functionality
	of the class
	\item private variables: variables that can only be accessed from within the class, they are used to store the state of
	the object and other internal data
	\item private functions: functions that can only be called from within the class, they are used to implement internal
	functionality of the class that is not exposed to the user, conventionally they are prefixed with an underscore to indicate
	that they are private
\end{itemize}

\newpage

\subsubsection*{Example of header file with class definition}
\begin{lstlisting}[language=C++]
	#ifndef CLASS_NAME_H // include guard to prevent multiple inclusions
	#define CLASS_NAME_H // of the header file

	#include <Arduino.h> // Arduino library with Arduino-specific functions and types

	class ClassName {
		public:
			ClassName(); // constructor
			void function1(); // public function or procedure
			int function2(int param); // public function with parameter and return value

		private:
			int variable1; // private variable
			float variable2; // private variable
			void privateFunction(); // private function
	};

	#endif // CLASS_NAME_H
\end{lstlisting}

\subsubsection*{Example of source file with class implementation}
\begin{lstlisting}[language=C++]
	#include <Arduino.h> // Arduino library with Arduino-specific functions and types
	#include "ClassName.h" // header file with class definition

	ClassName::ClassName() {
		// constructor implementation
	}

	void ClassName::function1() {
		// function implementation
	}

	int ClassName::function2(int param) {
		// function implementation
	}

	void ClassName::privateFunction() {
		// private function implementation
	}
\end{lstlisting}

\subsubsection*{Library file organisation and syntax highlighting}
To use a user-made library in an Arduino sketch, the library files (header and source files) must be placed in a specific folder
with the same name as the library in the Arduino libraries folder.

Usually, user-defined libraries contains symbols and variables that are not recognized by the Arduino IDE as keywords. To enable
syntax highlighting for those symbols, it is possible to add a keyword file in the library folder with the same name as the library
and the extension .txt that contains the list of keywords to be highlighted, for example:
\begin{lstlisting}
	--- inside ClassName.txt file ---
	ClassName    KEYWORD1
	function1    KEYWORD2
	function2    KEYWORD2
\end{lstlisting}

\newpage

\subsection{I\texorpdfstring{\textsuperscript{2}}{2}C Protocol}
\subsubsection*{Introduction to I\texorpdfstring{\textsuperscript{2}}{2}C Protocol}
The I\(^2\)C (Inter-Integrated Circuit) is a synchronous serial communication protocol that uses a master-slave architecture
where multiple devices can communicate over a common bus made of only two lines.

\subsubsection*{I\texorpdfstring{\textsuperscript{2}}{2}C connection lines}
The I\(^2\)C protocol uses two lines for communication:
\begin{itemize}
	\item SDA (Serial Data): line used to transmit data between devices
	\item SCL (Serial Clock): line used to synchronize the communication, is generated by the master device and can be forced
	low by the slave devices to pause the communication when they are not ready to receive or transmit data
\end{itemize}
NOTE: SDA line can change state only when SCL is LOW, exept for the start and stop conditions.

\subsubsection*{I\texorpdfstring{\textsuperscript{2}}{2}C communication process}
Every device on the I\(^2\)C bus has a unique 7-bit address that is used to identify the device during communication. If the
datasheet of the device specify an 8-bit address, a right shift of 1 bit must be applied (the least significant bit is eliminated).
The communication process works as follows:
\begin{itemize}
	\item the master device initiates the communication by sending a start condition (SCL keeps HIGH while SDA goes from
	HIGH to LOW) and starts generating the clock signal on the SCL line
	\item the master device sends the 7-bit address of the slave device it wants to communicate with, followed by a R/W bit
	(0 for write, 1 for read) to indicate if it wants to send or receive data from the slave and in the end the master
	releases the SDA line to allow the slave device to respond
	\item the slave device with the matching address responds with an ACK bit (0) if it is ready to communicate, otherwise
	it responds with a NACK bit (1) if it is not ready
	\item if the slave device ACKs the address, the master and the slave can exchange data bytes, according to the R/W bit
	sent by the master; after each byte of data, the receiver must send an ACK bit (0) or a NACK bit (1) to indicate whether
	transmission was successful or not
	\item after the communication is finished, the master sends a stop condition (SCL keeps HIGH while SDA goes from LOW to HIGH)
	and stops generating the clock signal on the SCL line
\end{itemize}

\subsubsection*{\texorpdfstring{\texttt{Wire.h}}{Wire} library to handle I\texorpdfstring{\textsuperscript{2}}{2}C communications}
In Arduino, the I\(^2\)C protocol is implemented through the \texttt{Wire} library, which provides functions to easily manage
I\(^2\)C communication. The library provides the following functions:
\begin{itemize}
	\item \texttt{Wire.begin()}: initializes the I\(^2\)C communication, must be called in the \texttt{setup()} function before
	using any other function of the library;
	\item \texttt{Wire.beginTransmission(address)}: initiates a transmission to the slave device with the specified address
	\item \texttt{Wire.write(data)}: enqueues a byte of data to the buffer for the slave device
	\item \texttt{Wire.endTransmission()} sends the bytes equeued in the buffer and ends the transmission
	\item \texttt{Wire.requestFrom(address, quantity)}: requests a specified quantity of bytes from the slave device with the
	specified address; returns the number of bytes received from the slave device
	\item \texttt{Wire.read()}: reads a byte of data received from the slave device; returns -1 if no data is available
\end{itemize}
The I\(^2\)C buffer is limited to 32 bytes, so if more data needs to be sent or received, it is necessary to split the data
into multiple transmissions.

There can be multiple Wire buses on the same Arduino board, for example on the Arduino Nano 33 BLE Sense Rev2 there is a second
I\(^2\)C bus connected by \texttt{Wire1.h} library, specifically designed to manage the built-in sensors on the board.

\subsubsection*{Using DS1621 digital temperature sensor with I\texorpdfstring{\textsuperscript{2}}{2}C communication}
The DS1621 is a digital temperature sensor that communicates over the I\(^2\)C bus. It has a 7-bit address that can be partially
configured using external pins of the sensor obtaining the 7-bit address in the following format: ``\texttt{1 0 0 1 A2 A1 A0}''.
I\(^2\)C commands for the DS1621 sensor are as follows:
\begin{itemize}
	\item \texttt{0xAC}: write the next byte into the configuration register of the sensor
	\item \texttt{0x02}: enable the continuous conversion mode of the sensor (temperature is continuously measured and updated in the sensor registers)
	\item \texttt{0xEE}: start a single temperature conversion (temperature is measured once and updated in the sensor registers)
	\item \texttt{0xAA}: read the temperature value from the sensor registers
\end{itemize}

Temperature values are represented in two bytes: the first byte contains the integer part of the temperature in two's complement
format, while the second byte contains the fractional part of the temperature in increments of 0.5\(^\circ C\) (0 for +0\(^\circ C\)
and 128 for +0.5\(^\circ C\)).

The following sketch shows how to use the \texttt{Wire} library to read the temperature value from the DS1621 sensor and print
it to the serial monitor.

\begin{lstlisting}[language=c++]
	#include <Wire.h> // I2C communication library

	#define DS1621_ADDR 0x48 // I2C address of the DS1621 sensor when A2, A1, A0 are LOW
	#define DS1621_CONFIG_REG 0xAC          // configuration register address
	#define DS1621_CONTINUOUS_CONV_CMD 0x02 // enable continuous conversion mode
	#define DS1621_START_CONV_CMD 0xEE      // start temperature conversion
	#define DS1621_READ_TEMP_CMD 0xAA       // read temperature value

	void setup() {
		Wire.begin(); // initialize I2C communication

		Wire.beginTransmission(DS1621_ADDR);    // start transmission to the sensor
		Wire.write(DS1621_CONFIG_REG);          // write to configuration register
		Wire.write(DS1621_CONTINUOUS_CONV_CMD); // enable continuous conversion mode
		Wire.endTransmission();                 // end transmission

		Wire.beginTransmission(DS1621_ADDR); // start transmission to the sensor
		Wire.write(DS1621_START_CONV_CMD);   // start temperature conversions
		Wire.endTransmission();              // end transmission

		Serial.begin(9600); // initialize serial communication for debugging
	}

	void loop() {
		Wire.beginTransmission(DS1621_ADDR); // start transmission to the sensor
		Wire.write(DS1621_READ_TEMP_CMD);    // send command to read temperature
		Wire.endTransmission();              // end transmission

		Wire.requestFrom(DS1621_ADDR, 2);    // request 2 bytes of data from the sensor

		byte byte1 = Wire.read(); // read the integer part of the temperature
		byte byte2 = Wire.read(); // read the fractional part of the temperature

		// compute the temperature value
		float temp = byte1 + (byte2 == 128 ? 0.5 : 0); 
		
		// print temperature value to serial monitor
		Serial.println("Temperature: " + temp + " C"); 

		delay(1000); // wait before reading the temperature again
	}
\end{lstlisting}

\newpage

\subsubsection*{Using TSL2561 light sensor with I\texorpdfstring{\textsuperscript{2}}{2}C communication}
The TSL2561 is a digital light sensor that communicates over the I\(^2\)C bus. It is equipped with two photodiodes on the same
CMOS chip that measure the visible and infrared light intensity with 16-bit resolution. There are two indipendent ADC channels
(0 and 1) to store respectively the visible light intensity and the infrared light intensity.

Values in the two channels registers are updated every 402ms. It is possible to further customize the integration time and the
gain of the sensor with additional commands (see datasheet).

The sensor has a 7-bit address that can be partially configured using an external pin of the sensor (0x29 when LOW,
0x39 when floating, 0x49 when HIGH). I\(^2\)C commands for the TSL2561 sensor are:
\begin{itemize}[itemsep=-1pt]
	\item \texttt{0x00}: write the next byte into the command register of the sensor
	\item \texttt{0x03}: power on the sensor (written on the command register)
	\item \texttt{0x00}: power off the sensor (written on the command register)
	\item \texttt{0x81}: access timing register of the sensor
	\item \texttt{0x02}: set low gain (1x) with 402ms integration time (written on the timing register)
	\item \texttt{0x8C}, \texttt{0x8D}: read channel 0 LSB (least significant byte) and MSB (most significant byte)
	\item \texttt{0x8E}, \texttt{0x8F}: read channel 1 LSB (least significant byte) and MSB (most significant byte)
\end{itemize}

The following sketch shows how to use the \texttt{Wire} library to read the light intensity values from the TSL2561 sensor and print
it to the serial monitor.

\begin{lstlisting}
	#include <Wire.h> // I2C communication library

	#define TSL2561_ADDR 0x29 // I2C address of the TSL2561 sensor when ADDR_PIN is LOW
	#define TSL2561_CMD_REG 0x00        // command register address
	#define TSL2561_POWER_ON 0x03       // power on the sensor
	#define TSL2561_POWER_OFF 0x00      // power off the sensor
	#define TSL2561_TIMING_REG 0x81     // timing register address
	#define TSL2561_LOW_GAIN_402MS 0x02 // set low gain with 402ms integration time
	#define TSL2561_READ_CH0_LSB 0x8C   // read channel 0 LSB
	#define TSL2561_READ_CH0_MSB 0x8D   // read channel 0 MSB
	#define TSL2561_READ_CH1_LSB 0x8E   // read channel 1 LSB
	#define TSL2561_READ_CH1_MSB 0x8F   // read channel 1 MSB

	void setup() {
		Wire.begin(); // initialize I2C communication

		Wire.beginTransmission(TSL2561_ADDR); // start transmission to the sensor
		Wire.write(TSL2561_CMD_REG);          // write to command register
		Wire.write(TSL2561_POWER_ON);         // power on the sensor
		Wire.endTransmission();               // end transmission

		Wire.beginTransmission(TSL2561_ADDR); // start transmission to the sensor
		Wire.write(TSL2561_TIMING_REG);       // write to timing register
		Wire.write(TSL2561_TIMING_LOW_GAIN);  // set low gain with 402ms integration time
		Wire.endTransmission();               // end transmission

		Serial.begin(9600); // initialize serial communication for debugging
	}

	void loop() {
		Wire.beginTransmission(TSL2561_ADDR); // start transmission to the sensor
		Wire.write(TSL2561_READ_CH0_LSB);     // send command to read channel 0 LSB
		Wire.endTransmission();               // end transmission
		Wire.requestFrom(TSL2561_ADDR, 1);    // request 1 bytes of data from the sensor
		byte ch0_lsb = Wire.read();           // read channel 0 LSB

		Wire.beginTransmission(TSL2561_ADDR); // start transmission to the sensor
		Wire.write(TSL2561_READ_CH0_MSB);     // send command to read channel 0 MSB
		Wire.endTransmission();               // end transmission
		Wire.requestFrom(TSL2561_ADDR, 1);    // request 1 bytes of data from the sensor
		byte ch0_msb = Wire.read();           // read channel 0 MSB
		
		float lum = 256* ch0_msb + ch0_lsb;   // compute lux value from channel values
		Serial.println("Lux: " + lum + " lx"); // print lux value to serial monitor
		delay(1000); // wait before reading the light intensity again
	}
\end{lstlisting}

\newpage

\subsection{Adafruit HX8357 LCD TFT display}
Adafruit HX8357 is an LCD TFT display with a resolution of 480x320 pixels with touchscreen capabilities. It can be controlled
via SPI communication and there is an Arduino library provided by Adafruit that allows to easily control the display.

\subsubsection*{Connections}
To connect the display to an Arduino board via SPI, the following connections must be made:
\begin{itemize}
	\item SCLK (Serial Clock) to the SCLK pin of the Arduino
	\item CIPO (Controller In Peripheral Out) to the MISO pin of the Arduino
	\item COPI (Controller Out Peripheral In) to the MOSI pin of the Arduino
	\item CS (Chip Select) to a digital pin of the Arduino (for example pin 10)
	\item DC (Data/Command) to a digital pin of the Arduino (for example pin 9); data/command pin is used to distinguish
	if the message transmitted to the display is data or command
	\item RST (Reset) to a digital pin of the Arduino (for example pin 8)
\end{itemize}

\subsubsection*{Object initialization}
To use the display functions in an Arduino sketch, it is necessary to import the library adn its dependencies and instantiate
an object of the class that manages the display, as in the following example:
\begin{lstlisting}
	#include <SPI.h>             // SPI protocol communication library
	#include <Adafruit_GFX.h>    // core graphics library for displays
	#include <Adafruit_HX8357.h> // library for the HX8357 display
	
	#define TFT_CS 10  // chip select pin
	#define TFT_DC 9   // data/Command pin
	#define TFT_RST 8  // reset pin
	
	// create an instance object of the display class
	Adafruit_HX8357 tft = Adafruit_HX8357(TFT_CS, TFT_DC, TFT_RST);

\end{lstlisting}

\subsubsection*{Basic library functions}
\begin{itemize}
	\item \texttt{tft.begin()}: initializes the display and must be called in the \texttt{setup()} function before using
	any other function of the library
	\item \texttt{tft.setRotation(i)}: sets the rotation of the display, where \(i\) can be in range 0-3 and corresponds
	to a specific orientation of the display (0°, 90°, 180°, 270°)
	\item \texttt{tft.fillScreen(color)}: fills the entire display with the specified color
	\item \texttt{tft.setTextColor(color)}: sets the color for text rendering
	\item \texttt{tft.setTextSize(size)}: sets the size for text rendering, where \(size\) is an integer that scales the
	default text size (1 is default size, 2 is double size, etc.); default text size is 5x8 pixels.
	\item \texttt{tft.setCursor(x, y)}: sets the cursor position for text rendering at the specified coordinates
	\item \texttt{tft.print(text)}: prints the text at the current cursor position with current text, color and size
	\item \texttt{tft.println(text)}: same as \texttt{tft.print()}, plus moves the cursor to the beginning of next line
	\item \texttt{tft.drawPixel(x, y, color)}: draws a pixel at the specified coordinates with the specified color
	\item \texttt{tft.drawLine(x0, y0, x1, y1, color)}: draws a line based on the specified parameters
	\item \texttt{tft.drawRect(x, y, w, h, color)}: draws a rectangle based on the specified parameters
	\item \texttt{tft.fillRect(x, y, w, h, color)}: draws a filled rectangle based on the specified parameters
	\item \texttt{tft.drawCircle(x, y, r, color)}: draws a circle based on the specified parameters
	\item \texttt{tft.fillCircle(x, y, r, color)}: draws a filled circle based on the specified parameters
\end{itemize}

\subsubsection*{Pixel organization}
The display is organized in a grid of pixels, where the top-left corner is the origin (0, 0) and the bottom-right corner is
(479, 319). Every text character is rendered from the current cursor pixel-position to bottom-right corner of the display.

\subsubsection*{Color coding}
The display is capable of 18-bit color depth, but the color is usually specified in 16-bit format (5 bits for red, 6 bits for
green and 5 bits for blue) to save memory and bandwidth. Standard colors are pre-defined in the library as constants, for example:
\begin{center}
	\begin{minipage}{0.45\textwidth}
		\begin{itemize}
			\item \texttt{HX8357\_BLACK}: black color (0x0000)
			\item \texttt{HX8357\_RED}: red color (0xF800)
			\item \texttt{HX8357\_GREEN}: green color (0x07E0)
			\item \texttt{HX8357\_ BLUE}: blue color (0x001F)
		\end{itemize}
	\end{minipage}
	\begin{minipage}{0.5\textwidth}
		\begin{itemize}
			\item \texttt{HX8357\_CYAN}: cyan color (0x07FF)
			\item \texttt{HX8357\_MAGENTA}: magenta color (0xF81F)
			\item \texttt{HX8357\_YELLOW}: yellow color (0xFFE0)
			\item \texttt{HX8357\_WHITE}: white color (0xFFFF)
		\end{itemize}
	\end{minipage}
\end{center}
It is possile to create custom colors with the \texttt{tft.color565(r, g, b)} library function, where \(r\), \(g\) and \(b\)
are the red, green and blue components of the color in range (\([0, 255]\)).

\subsubsection*{Touchscreen functions}
The display has a built-in resistive touchscreen that can be used to detect touch events and get the coordinates of the touch.
Touchscreen is realized by two layers of conductive material separated by a small gap. When the screen is touched, the two layers
come into contact and the resistance between them changes, allowing to get the coordinates and the pressure of the touch.

The LCD has 4 pins for the touchscreen (XP, XM, YP, YM) that can be connected to the Arduino (D7, A3, A2, D8 respectively)
and managed with the following library functions:
\begin{itemize}
	\item \texttt{Touchscreen(XP, YP, XM, YM, resistance)}: constructor of the touchscreen class that initializes the touchscreen
	object with the specified pins and the resistance value of the touchscreen
	\item \texttt{ts.getPoint()}: function that returns a structure with the coordinates (x, y, pressure) in range (\([0, 4095]\))
	for the x and y coordinates and in range (\([0, 255]\)) for the pressure value; if the pressure value is greater than the
	\texttt{ts.pressureThreshhold} (equal to 0), the touch is considered valid
\end{itemize}

The following code snippet shows how to get the coordinates of a touch event and map them to the display coordinates.
\begin{lstlisting}
	Touchscreen ts = Touchscreen(XP, YP, XM, YM, resistance); // object constructor
	TSPoint p = ts.getPoint(); // get the coordinates of the touch (x, y, pressure)
	p.x = map(p.x, 0, 4095, 0, tft.width()); // map x coordinate to display width
	p.y = map(p.y, 0, 4095, 0, tft.height()); // map y coordinate to display height

	// check if a touch event occured by comparing the pressure value to the threshold
	if (p.pressure > ts.pressureThreshhold) {
		// touch event has occured at coordinates (p.x, p.y) with p.pressure
	}
\end{lstlisting}

\newpage
\footnotesize

\subsection{Heart rate monitor}
\subsubsection*{Introduction to heart rate monitoring on Arduino}
The idea for building a heart rate monitor is to detect the heart rate of a person by measuring the changes in reflectance of red and
infrared light on the skin caused by the blood flow in the arteries. To do this an optical sensor is used in pair with a red LED and
an infrared LED. The photodiode signal has a small pulsation component superimposed on a large DC component, so the signal must be
conditioned and filtered to extract the pulsation component and get a clean signal that can be read by the ADC of the Arduino board. 

\subsubsection*{Circuit components}
\begin{center}
	\begin{tikzpicture}[scale=0.7, transform shape]
		\draw (5.5, 8) to[american resistor, *-, l={$R_{IR}$}] (5.5, 6);
		\node[ground] at (5.5, 4.25){};
		\draw (5.5, 6) to[empty led] (5.5, 4);
		\draw (7.75, 8) to[american resistor, *-, l={$R_{RED}$}] (7.75, 6);
		\node[ground] at (7.75, 4.25){};
		\draw (7.75, 6) to[empty led] (7.75, 4);
		\node[shape=rectangle, minimum width=2.465cm, minimum height=0.715cm] at (5.5, 8.375){} node[anchor=north, align=center, text width=2.077cm, inner sep=6pt] at (5.5, 8.75){\small IR LED pin};
		\node[shape=rectangle, minimum width=2.465cm, minimum height=0.715cm] at (7.75, 8.375){} node[anchor=north, align=center, text width=2.077cm, inner sep=6pt] at (7.75, 8.75){\small RED LED pin};
		\draw (10, 4) to[empty photodiode, i^<=$I_{PD}$, current/distance=0.6] (10, 6);
		\node[ground] at (10, 4.25){};
		\node[op amp] at (12.06, 6){\small opamp 1};
		\node[ground] at (10.87, 5.51){};
		\draw (10, 6) |- (10.87, 6.49);
		\draw (10.87, 7.99) to[variable american resistor, l={$R_{pot}$}] (13.25, 8);
		\draw (13.25, 6) to[capacitor, l={$C_1$}] (14.75, 6);
		\draw (14.75, 6) to[american resistor, -*, l_={$R_4$}] (14.75, 8);
		\draw (14.75, 6) to[american resistor, l_={$R_1$}] (14.75, 4);
		\node[op amp, yscale=-1] at (16.69, 5.51){} node[anchor=center] at (16.58, 5.51) {\small opamp 2};
		\node[ground] at (14.75, 4.25){};
		\draw (10.87, 7.99) -| (10.87, 6.49);
		\draw (13.25, 8) -| (13.25, 6);
		\draw (14.75, 6) -- (15.5, 6);
		\node[ground] at (11.977, 5.461){};
		\draw (11.977, 6.539) to[short, -*] (11.977, 7.039);
		\node[shape=rectangle, minimum width=1.215cm, minimum height=0.426cm] at (11.977, 7.269){} node[anchor=north, align=center, text width=1.109cm, inner sep=2pt] at (11.977, 7.5){\small $V_{CC}$};
		\node[shape=rectangle, minimum width=1.215cm, minimum height=0.426cm] at (14.75, 8.23){} node[anchor=north, align=center, text width=1.109cm, inner sep=2pt] at (14.75, 8.461){\small $V_{CC}$};
		\node[ground] at (16.607, 4.971){};
		\draw (16.607, 6.049) to[short, -*] (16.607, 6.549);
		\node[shape=rectangle, minimum width=1.215cm, minimum height=0.426cm] at (16.607, 6.779){} node[anchor=north, align=center, text width=1.109cm, inner sep=2pt] at (16.607, 7.01){\small $V_{CC}$};
		\draw (18, 5.5) to[american resistor, l_={$R_2$}] (18, 3.5);
		\draw (18.75, 5.5) to[capacitor, l={$C_2$}] (18.75, 3.5);
		\draw (15.5, 3.5) to[american resistor, l={$R_3$}] (15.5, 1.5);
		\node[ground] at (15.5, 1.75){};
		\draw (15.5, 5.02) -- (15.5, 3.5);
		\draw (18.75, 3.5) -- (15.5, 3.5);
		\draw (19.5, 5.5) to[capacitor, l={$C_1$}] (21, 5.5);
		\draw (21, 5.5) to[american resistor, -*, l_={$R_4$}] (21, 7.5);
		\draw (21, 5.5) to[american resistor, l_={$R_1$}] (21, 3.5);
		\node[op amp, yscale=-1] at (22.94, 5.01){} node[anchor=center] at (22.83, 5.01) {\small opamp 3};
		\node[ground] at (21, 3.75){};
		\draw (21, 5.5) -- (21.75, 5.5);
		\node[shape=rectangle, minimum width=1.215cm, minimum height=0.426cm] at (20.75, 7.73){} node[anchor=north, align=center, text width=1.109cm, inner sep=2pt] at (20.75, 7.961){\small $V_{CC}$};
		\node[ground] at (22.857, 4.471){};
		\draw (22.857, 5.549) to[short, -*] (22.857, 6.049);
		\node[shape=rectangle, minimum width=1.215cm, minimum height=0.426cm] at (22.857, 6.279){} node[anchor=north, align=center, text width=1.109cm, inner sep=2pt] at (22.857, 6.51){\small $V_{CC}$};
		\draw (24.25, 5) to[american resistor, l_={$R_2$}] (24.25, 3);
		\draw (25, 5) to[capacitor, l={$C_2$}] (25, 3);
		\draw (21.75, 3) to[american resistor, l={$R_3$}] (21.75, 1);
		\node[ground] at (21.75, 1.25){};
		\draw (21.75, 4.52) -- (21.75, 3);
		\draw (25, 3) -- (21.75, 3);
		\draw (17.88, 5.51) to[short] (18, 5.5);
		\draw (18, 5.5) -- (19.5, 5.5);
		\draw (24.13, 5.01) to[short] (24.25, 5);
		\draw (24.25, 5) to[short, -*] (25.75, 5);
		\node[shape=rectangle, minimum width=1.465cm, minimum height=0.83cm] at (25.75, 5.5){} node[anchor=north, align=center, text width=1.359cm, inner sep=2pt] at (25.75, 5.933){\small Analog input pin};
		\node[shape=rectangle, minimum width=1.09cm, minimum height=1.09cm] at (7, 4.687){} node[anchor=north, align=center, text width=0.702cm, inner sep=6pt] at (7, 5.25){\small RED LED};
		\node[shape=rectangle, minimum width=1.09cm, minimum height=1.09cm] at (4.75, 4.687){} node[anchor=north, align=center, text width=0.702cm, inner sep=6pt] at (4.75, 5.25){\small IR LED};
	\end{tikzpicture}
\end{center}

\noindent
Circuit is composed by:
\begin{enumerate}
	\item two LEDs (red and infrared) connected in series with two resistors to limit the current through them, controlled by two digital
	pins of the Arduino board to turn them on and off
	\item a photodiode connected in reverse bias to a transimpedance amplifier (opamp 1) with a potentiometer in the feedback loop that
	converts the photodiode current into a voltage signal
	\item a passive high pass filter with a cutoff frequency of \(f_C \approx 0.8 \; Hz\) and a voltage offset of \(0.15 \; V\) (\(C_1\),
	\(R_1\) and \(R_4\)) to remove the DC component of the signal and apply a small positive offset to the signal to make it suitable
	for op amps with single positive power supply
	\item an active low pass filter with a cutoff frequency of \(f_C \approx 3 \; Hz\) and a gain of \(A_v = 11\) (opamp 2, \(R_2\), \(R_3\)
	and \(C_2\)) to filter the signal, remove high frequency noise and amplify the pulsation component of the signal
	\item a second passive high pass filter with offset (\(C_1\), \(R_1\) and \(R_4\))
	\item a second active low pass filter (opamp 3, \(R_2\), \(R_3\) and \(C_2\))
	\item the output terminal connected to an analog input pin of the Arduino board
\end{enumerate}

\subsubsection*{Oxygen saturation monitoring}
Oxygen saturation can be estimated by measuring the ratio of the pulsation component of the red and infrared signals, since oxygenated
and deoxygenated hemoglobin have different light absorption coefficients for red and infrared light. The ratio can be computed with the
following formula:
\[R = \frac{\displaystyle \left(\frac{AC}{DC}\right)_{RED}}{\displaystyle \left(\frac{AC}{DC}\right)_{IR}} = \frac{\displaystyle \frac{V_{\text{max, RED}} - V_{\text{min, RED}}}{V_{\text{min, RED}}}}{\displaystyle \frac{V_{\text{max, IR}} - V_{\text{min, IR}}}{V_{\text{min, IR}}}} \qquad\qquad \text{SpO}_2 = 97.94 + 1.15 \cdot R\]